\textcolor{black}{\section{Introducción}}

\noindent 

La constante de Planck, $h$, es una de las constantes fundamentales de la física moderna, ya que establece la cuantización de la energía electromagnética. En los fenómenos microscópicos, la energía de un fotón se relaciona con su frecuencia mediante la expresión:

\begin{equation}
	E=h \nu,
	\label{eq:planck}
\end{equation}

donde $\nu$ es la frecuencia de la radiación emitida o absorbida.

Por otro lado, un diodo LED (\textit{Light Emitting Diode}), son dispositivos semiconductores formados por la unión de un material tipo $P$, rico en huecos como portadores mayoritarios, y un material tipo $N$, en el que predominan los electrones libres. Al polarizar directamente la unión, los electrones y huecos son inyectados hacia la región de agotamiento, donde pueden recombinarse. En un diodo LED, esta recombinación entre los huecos y los electrones libera energía en forma de fotones, dando lugar a la emisión de luz.

Para que la emisión de luz ocurra es necesario aplicar una diferencia de potencial mínima conocida como \textit{voltaje umbral}. Este voltaje corresponde aproximadamente a la energía que deben adquirir los portadores para atravesar la barrera de potencial de la unión $P-N$ y poder recombinarse. En consecuencia, la energía eléctrica suministrada a cada electrón puede escribirse como:

\begin{equation}
	E = eV_U
\end{equation}
donde $V_U$ es el voltaje umbral del LED.

Usando la ecuación \eqref{eq:planck} se llega a la relación entre el voltaje y la frecuencia de la luz emitida.

\begin{equation}
	V_U =  \frac{h}{e} \nu
\end{equation}

La frecuencia de la luz emitida puede obtenerse a partir de su longitud de onda utilizando la relación

\begin{equation}
	f=\frac{c}{\lambda},
\end{equation}

donde $c$ es la velocidad de la luz en el vacío.